\documentclass{article}
\usepackage{amsmath,amssymb,hyperref}
\title{Shared boundaries derived from independent connected lifts}
\date{15 November 2026}
\begin{document}\maketitle
\paragraph{Scope.} A written binary theorem with finite shared-product closure. It derives the comma-connectivity condition needed by selective right-Kan assembly from separate component conditions. It is an ordinary categorical theorem, including normal lax Prof presentations with possibly nonrepresentable transitions. It is not an automatic extension to arbitrary bicategorical enrichment, unrestricted coupled processes, or feedback.

\paragraph{Independent data.} Let $D$ be a small nonempty sifted category in the classical sense: for every $b,c$, the category of cospans $b\to d\leftarrow c$, with commuting cospan maps, is nonempty connected. Let $p:E\to D$ and $q:F\to D$ be functors. For $e\in E$ and $a:p(e)\to d$, define $L_p(e,a)$ to have objects arrows $u:e\to e'$ over $a$, and morphisms vertical arrows $v:e'\to e''$ satisfying $vu=u'$. Require every such category to be nonempty connected, and likewise for $q$. These are conditions on each component separately. They do not prescribe a joint coupling or assume the final shared comparison is exact.

\paragraph{Theorem.} Let $S=E\times_D F$ be the strict shared-boundary category, and $i:S\to E\times F$ its inclusion. For every $(e,f)$, the comma category $(e,f)\downarrow i$ is nonempty connected.

\paragraph{Proof.} Send an object $(u:e\to e',v:f\to f')$ of that comma category to the boundary cospan
\[
 p(e)\xrightarrow{p(u)}d\xleftarrow{q(v)}q(f),\qquad d=p(e')=q(f').
\]
The fiber over a cospan $(a,b)$ is exactly $L_p(e,a)\times L_q(f,b)$, hence nonempty connected. Every cospan map $h:d\to d'$ lifts out of every object in that fiber: choose an $h$-lift from $e'$ and from $f'$, whose composites with $u,v$ form an object over the new cospan. The two chosen arrows give a morphism of the total comma category. No uniqueness or opcartesian universal property is required.

The base cospan category is nonempty connected by siftedness. Consequently the total category is nonempty. To connect any two total objects, take a finite zigzag between their base cospans. Each oriented edge has an outgoing lift from any chosen object over its source. When traversing an edge backwards, choose a lift from an object over the preceding source and connect its target to the current object inside the connected target fiber. Forward steps and fiber zigzags are treated similarly. This produces a zigzag between the two original total objects and proves connectedness.

\paragraph{Consequences for the selected comparison adjoint.} In the ordinary module action realization, the conjoint $i^*$ acts by restriction on covariant presheaves. Its selected right adjoint is $\operatorname{Ran}_i$. The theorem therefore establishes the connected-comma condition used to commute this right Kan extension with independent separated external-product assembly. It does not claim that $i^*$ already has a right adjoint in Prof, nor that $\operatorname{Ran}_i$ is cocontinuous. The argument supplies the missing joint condition from local component hypotheses rather than assuming it.

\paragraph{Closure under finite shared products.} The shared projection $s:E\times_D F\to D$ again has the lift property: for an object $(e,f)$ and base arrow $a$, its lift category is $L_p(e,a)\times L_q(f,a)$. Thus binary shared assembly stays in the admitted class. For any finite nonempty family, iterated shared products are canonically equivalent by tuple reassociation. Their inclusion into the independent product factors into binary shared inclusions and products with identities. Connected-comma functors are closed under products and composition, as also follows from their canonical preservation of constant presheaves by right Kan extension. Hence each such finite inclusion has connected nonempty commas. Right Kan extensions along the different factorizations agree coherently because they are right adjoints to the same reassociated restriction functor. This is coherence of these selected restriction adjoints, not a claim of a total tensor on all newly admitted arrows.

\paragraph{Normal lax module presentations.} For a normal lax functor $H:D\to\mathrm{Prof}$, the enriched cone/Grothendieck category has objects $(d,x)$ and arrows over $a:d\to d'$ given by elements of $H(a)(-,x)$. The lift category is the category of elements of that covariant transition weight, with target-fiber action. Thus the theorem applies when those categories of elements are nonempty connected. Representable transition weights satisfy this, but representability is not necessary. For example take $D=(0\to1)$, $H0=1$, $H1=BG$ for a nontrivial group, and transition weight the singleton trivial $G$-set. Its category of elements is $BG$, connected nonempty but not representable. Identity-arrow lift categories have initial identity lifts. This diagnostic witnesses genuine scope beyond functorial/opcartesian transport; it is not the theorem's premise.

\paragraph{Source and boundary.} The siftedness criterion is classical Ad\'amek--Rosick\'y background; the distinction from homotopy siftedness is stated explicitly in Kerodon, Warning 9.7.2.2, \url{https://kerodon.net/tag/02QD}. Only classical connectivity is used here. The total-connectivity proof is given in full and no novelty claim is made. The result does not assert all normal lax Prof diagrams have connected lifts, does not preserve arbitrary hostile choices or quantitative laws, and does not construct a feedback operation. A general enriched analogue and the interaction with the actual original admitted pure interface remain separate obligations.
\end{document}
